\documentclass[11pt]{article}
\usepackage[margin=1in]{geometry}
\usepackage{amsmath,amssymb,amsthm}
\usepackage[hidelinks]{hyperref}

\newtheorem{theorem}{Theorem}
\newtheorem{lemma}{Lemma}
\newtheorem{corollary}{Corollary}

\newcommand{\N}{\mathbb{N}}
\newcommand{\ellone}{\ell_1}
\newcommand{\NN}{\N^\N}
\newcommand{\co}{\operatorname{co}}

\title{The free locally convex space \(L(\mathbb N^\mathbb N)\) is not multi-Rosenthal}
\author{arXiv automatic math run \texttt{fa\_banach\_001}}
\date{2026-06-15}

\begin{document}
\maketitle

\section*{Source problem}

The source paper is M. Komisarchik and M. Megrelishvili,
\emph{Tameness and Rosenthal type locally convex spaces},
arXiv:2203.02368.  On page 33, Question 8.7 asks:

\begin{quote}
Examine if \(L(\mathbb N^{\mathbb N})\) is multi-Asplund or, at least, multi-Rosenthal.
\end{quote}

Here \(L(X)\) denotes the free locally convex space over the Tychonoff space
\(X\), and \(\mathbb N^{\mathbb N}\) is the Baire space, homeomorphic to the
space of irrationals \(P\) mentioned immediately before the question.

\section*{Definitions}

A Banach space is called \emph{Rosenthal} if it contains no isomorphic copy of
\(\ell_1\).  A locally convex space \(E\) is \emph{multi-Rosenthal} if it is
isomorphic, as a locally convex space, to a subspace of a product of Rosenthal
Banach spaces.  Similarly, \(E\) is \emph{multi-Asplund} if it embeds into a
product of Asplund Banach spaces.

Every Asplund Banach space is Rosenthal: if it contained \(\ell_1\), then the
separable subspace \(\ell_1\) would have nonseparable dual \(\ell_\infty\),
contradicting the separable-subspace characterization of Asplund spaces.
Therefore multi-Asplund implies multi-Rosenthal.

\section*{Idea of the Proof}

The proof uses the universal property of \(L(X)\) in the most direct possible
way.  If \(L(\NN)\) were multi-Rosenthal, then every Banach-valued continuous
map on \(\NN\) would extend to a continuous linear operator on \(L(\NN)\), and
that operator would be controlled by finitely many coordinates of the product
of Rosenthal spaces.  Thus the operator would factor through a Rosenthal Banach
space.

But the Baire space maps continuously onto every Polish space, in particular
onto \(\ell_1\).  Extending such a surjection would force \(\ell_1\) to be a
quotient of a Rosenthal Banach space.  This is impossible, because the property
of containing no \(\ell_1\) passes to Banach quotients.

\section*{Auxiliary facts}

\begin{lemma}[finite-coordinate factorization]
\label{lem:factor}
Let \(E\) be a multi-Rosenthal locally convex space and let
\(T:E\to B\) be a continuous linear map into a Banach space.  Then \(T\)
factors through a Rosenthal Banach space.
\end{lemma}

\begin{proof}
Embed \(E\) as a locally convex subspace of a product
\(\prod_{i\in I} R_i\), where each \(R_i\) is a Rosenthal Banach space.  Since
\(T\) is continuous into a normed space, there are finitely many coordinates
\(F\subset I\) and a constant \(C\) such that
\[
 \|Tx\|\le C\max_{i\in F}\|x_i\|_{R_i}
\]
for every \(x\in E\).  Hence \(T\) vanishes on the kernel of the coordinate
projection \(\pi_F:E\to \prod_{i\in F}R_i\), and so \(T=S\pi_F\) for a bounded
linear map \(S\) on \(\pi_F(E)\).  Extending \(S\) to the closure
\[
G=\overline{\pi_F(E)}\subseteq \prod_{i\in F}R_i
\]
gives a factorization \(E\to G\to B\).

The finite product of Rosenthal Banach spaces is Rosenthal, and closed
subspaces of Rosenthal Banach spaces are Rosenthal.  Thus \(G\) is Rosenthal.
\end{proof}

\begin{lemma}[quotients]
\label{lem:quotient}
A Banach quotient of a Rosenthal Banach space is Rosenthal.
\end{lemma}

\begin{proof}
Let \(q:X\to Y\) be a quotient map, and suppose that \(Y\) contains an
isomorphic copy \(J:\ell_1\to Y\).  By the open mapping theorem, there is
\(C>0\) such that for every \(y\in Y\) one can choose \(x\in X\) with
\(qx=y\) and \(\|x\|\le C\|y\|\).  Choose \(x_n\in X\) with
\(qx_n=J e_n\) and \(\sup_n\|x_n\|<\infty\).

Define \(S:\ell_1\to X\) by
\[
S(a)=\sum_{n=1}^{\infty} a_n x_n .
\]
The series converges absolutely, so \(S\) is bounded, and \(qS=J\).  Since
\(J\) is an isomorphic embedding,
\[
\|Sa\|\ge \frac{1}{\|q\|}\|J a\|
\]
for every \(a\in\ell_1\).  Hence \(S\) embeds \(\ell_1\) into \(X\).  Thus if
\(X\) contains no copy of \(\ell_1\), neither can \(Y\).
\end{proof}

\section*{Main theorem}

\begin{theorem}
\label{thm:main}
The free locally convex space \(L(\mathbb N^{\mathbb N})\) is not
multi-Rosenthal.  Consequently it is not multi-Asplund.
\end{theorem}

\begin{proof}
The Baire space \(\NN\) admits a continuous surjection
\[
s:\NN\to\ell_1,
\]
because every Polish space is a continuous image of \(\NN\).

By the universal property of the free locally convex space, \(s\) extends to a
continuous linear operator
\[
\widehat{s}:L(\NN)\to\ell_1
\]
with \(\widehat{s}(\delta_x)=s(x)\).  Since \(s\) is onto, \(\widehat{s}\) is
onto.

Assume, toward a contradiction, that \(L(\NN)\) is multi-Rosenthal.  By
Lemma~\ref{lem:factor}, \(\widehat{s}\) factors as
\[
L(\NN)\longrightarrow G \stackrel{u}{\longrightarrow} \ell_1
\]
where \(G\) is a Rosenthal Banach space and \(u\) is bounded linear.  Since
\(\widehat{s}\) is onto, \(u(G)=\ell_1\).  Therefore \(\ell_1\) is a Banach
quotient of the Rosenthal Banach space \(G\), contradicting
Lemma~\ref{lem:quotient}.  Thus \(L(\NN)\) is not multi-Rosenthal.

Finally, every Asplund Banach space is Rosenthal, so a multi-Asplund embedding
would in particular be a multi-Rosenthal embedding.  Hence \(L(\NN)\) is not
multi-Asplund.
\end{proof}

\section*{Verification notes}

The proof uses only standard permanence properties and one classical
descriptive-set-theoretic fact: every Polish space is a continuous image of
the Baire space.  The quotient permanence of Rosenthal Banach spaces is proved
above using the projectivity of \(\ell_1\), unfolded into an explicit lifting
of its unit vector basis.

The result is stronger than the previously known statement quoted by
Komisarchik--Megrelishvili that \(L(P)\), with \(P\) the irrationals, is not
multi-reflexive.  It rules out embeddings into products of all Rosenthal
Banach spaces, not only reflexive Banach spaces.

\section*{Novelty and literature check}

Local run indexes were searched for arXiv:2203.02368, the source title,
Question 8.7, \(L(\mathbb N^\mathbb N)\), multi-Asplund, multi-Rosenthal, and
free locally convex space.  No existing packet or attempt for this question
was found.

Bounded web searches on 2026-06-15 for exact phrases including
\(\text{``}L(\mathbb N^\mathbb N)\text{''}\) with multi-Rosenthal,
multi-Asplund with free locally convex spaces, and the source title with
multi-Rosenthal found the source paper and nearby multi-reflexive/nuclear
literature, but no prior answer to Question 8.7.  The novelty confidence is
moderate: the argument is short and may be known to specialists, but no exact
published answer was located in this search.

\section*{References}

\begin{itemize}
\item M. Komisarchik and M. Megrelishvili, \emph{Tameness and Rosenthal type
locally convex spaces}, arXiv:2203.02368.
\item A. Leiderman and V. Uspenskij, \emph{Is the free locally convex space
\(L(X)\) nuclear?}, arXiv:2106.13413.
\item J. Diestel, \emph{Sequences and Series in Banach Spaces}, for standard
Rosenthal-space and \(\ell_1\) background.
\item A. S. Kechris, \emph{Classical Descriptive Set Theory}, for the
classical fact that every Polish space is a continuous image of the Baire
space.
\end{itemize}

\section*{Human review recommendation}

This packet should be reviewed as a likely full negative answer to Question
8.7.  The main points to check are Lemma~\ref{lem:factor}, especially the
finite-coordinate domination step for maps from a subspace of a product, and
Lemma~\ref{lem:quotient}, the quotient permanence of the Rosenthal property.

\end{document}
